% ECONOMICS_PROBLEM: {"id": "D91-530", "title": "A portable model of cognitive biases", "jel_code": "D91", "source_id": "OP-0301"}
% !TeX program = xelatex
\documentclass[11pt,a4paper]{article}
\usepackage[margin=23mm]{geometry}
\usepackage{fontspec}
\setmainfont{Latin Modern Roman}
\usepackage{amsmath,amssymb,mathtools}
\usepackage{xcolor,enumitem,booktabs,longtable,array,xurl,hyperref}
\definecolor{muted}{HTML}{53636C}
\hypersetup{colorlinks=true,urlcolor=blue,linkcolor=blue}
\setlength{\parindent}{0pt}
\setlength{\parskip}{5pt}
\newcommand{\R}{\mathbb R}
\newcommand{\E}{\mathbb E}
\newcommand{\N}{\mathbb N}
\newcommand{\ind}{\mathbf 1}
\newcommand{\Var}{\operatorname{Var}}
\newcommand{\Cov}{\operatorname{Cov}}
\newcommand{\supp}{\operatorname{supp}}
\DeclareMathOperator*{\argmin}{arg\,min}
\DeclareMathOperator*{\argmax}{arg\,max}
\newcommand{\entrymeta}[3]{{\sffamily\small\color{muted}\textbf{\texttt{#1}}\quad|\quad #2\par #3\par}}
\newcommand{\Problemref}[1]{\texttt{#1}}
\newcommand{\JELref}[2]{\texttt{#2} (JEL \texttt{#1})}
\begin{document}
\section*{A portable model of cognitive biases}
\textbf{Problem ID:} \texttt{D91-530}\quad \textbf{JEL:} \texttt{D91}\par
% BEGIN ECONOMICS STATEMENT
\label{problem:OP-0301}
{\sffamily\small\color{muted}Original subject group: Behavioral and experimental economics\par}
\label{definitions:problem:30007583}
\textbf{Audit classification:} Published research agenda. Enke–Graeber §8 explicitly discusses cognitive defaults and noise-measure validity. The cross-domain shrinkage and prediction test are editorial.
\subsection*{Definitions and precise research target}
\textbf{Model and research target.} Let \(d\) index risky choice, ambiguous choice, belief updating, and intertemporal choice. In each domain individual \(i\) faces a task with observed attributes \(X_{id}\), gives response \(Y_{id}\), and reports cognitive uncertainty \(Q_{id}\in[0,1]\), defined as subjective doubt that the chosen action maximizes expected utility. Let \(m_d(X_{id};\theta_i)\) denote a domain-specific benchmark decision, where \(\theta_i\) contains stable preference or belief parameters. A candidate shared mechanism has the form
\[Y_{id}=(1-\lambda_{id})m_d(X_{id};\theta_i)+\lambda_{id}b_d(X_{id};\eta)+\varepsilon_{id},\qquad \lambda_{id}=g(Q_{id},X_{id};\gamma).\]
Here \(b_d\) is a cognitive default, \(g\) maps measurements to shrinkage intensity, and \(\eta,\gamma\) are parameters shared across specified domains. Define prediction loss in a held-out domain \(d\) as \(R_d=\mathbb E[\ell(Y_{id},\widehat Y_{id})]\), for a preregistered loss \(\ell\); estimate the prediction using parameters trained without that domain. Can measured uncertainty and a low-dimensional account of defaults and task difficulty improve \(R_d\) over independently fitted domain models, while preserving calibration across populations and description-versus-experience treatments? Also estimate whether \(Q_{id}\) predicts independently measured response noise. This formulation asks for testable portability and calibration rather than asserting that a universal model exists. Utility benchmarks and scaling of responses must be specified separately before the cross-domain comparison has empirical meaning.

\emph{Definitions and maintained assumptions (editorial unless a source definition is named).} Restrict \(Y_{id}\) to scalar responses on a declared domain-specific scale, or replace the convex-combination equation by a model on a shared feasible vector space. \(m_d\) and \(b_d\) must lie in the same response space; for discrete choices a mixture of probabilities, not a fractional action, is appropriate. Require \(g(Q,X;\gamma)\in[0,1]\), integrable disturbance, and a stated conditional mean restriction if regression parameters are to be interpreted. Define measured \(Q\) by the exact incentivized or survey question and reporting-error kernel; it is not assumed to equal true cognitive noise. Stable individual parameters \(\theta_i\) require independent identification or prespecified pooling. Held-out-domain risk \(R_d\) uses training independent of test observations; domain-specific benchmark models must receive comparable training sample sizes and tuning budgets. Calibration is a declared conditional prediction condition or bin-based approximation. The substantive hypothesis is reduction in out-of-domain loss by a restricted mechanism, not the unconstrained ability of a flexible model to fit existing responses.
\subsection*{Variants and assumption changes}
\begin{enumerate}
\item \textbf{Fixed defaults (editorial).} Prespecify a default per domain and a shared finite-dimensional shrinkage function. Test whether it predicts held-out risk and belief tasks better than equally constrained baselines.
\item \textbf{Noise validation (source agenda, editorial design).} Obtain repeated independent responses and uncertainty reports on the same tasks. Identify when reported uncertainty predicts independently measured variance rather than only systematic central tendency.
\end{enumerate}
\subsection*{References and source audit}
\begin{enumerate}
\item §8 printed pp.32–34 in March2023 manuscript. Supports the stated research area; exact displayed specification is editorial. \url{https://thomasgraeber.com/assets/papers/EnkeGraeber_CognitiveUncertainty.pdf}
\end{enumerate}
\begin{enumerate}\raggedright
\item\hypertarget{definitions:source:30007583:1}{} Benjamin Enke; Thomas Graeber. 2023. \emph{Cognitive Uncertainty}. §8 printed pp.32–34 in March2023 manuscript.
\url{https://thomasgraeber.com/assets/papers/EnkeGraeber_CognitiveUncertainty.pdf}.
DOI: \url{https://doi.org/10.1093/qje/qjad025}.
\end{enumerate}

\subsection*{Common conventions: Source mathematical conventions}

These conventions apply to each entry unless explicitly replaced.

\textbf{Models and identification.} A primitive model \(m\) specifies all probability laws, feasible choices, information, equilibrium and measurement rules used by its target. The admissible class \(\mathcal M\) is fixed independently of outcomes. If \(P_O(m)\) is its observable probability law and \(\tau(m)\) the target, its identified set at observable law \(P\) is
\[
 \mathcal I_\tau(P)=\{\tau(m):m\in\mathcal M,\ P_O(m)=P\}.
\]
Point identification means that this set is a singleton. An empty set signals incompatibility of data and assumptions. Coordinate bounds need not describe the joint set. Bounds are sharp only if every retained target is attainable or arbitrarily approachable by compatible models. Finite bounds require explicit restrictions on unobserved quantities; unrestricted confounding, hidden wealth, extrapolation or arbitrary equilibrium selection can yield uninformative sets.

\textbf{Causal interventions.} A potential outcome \(Y_i(p)\) is the outcome under a complete policy regime \(p\), including timing, financing, information and market spillovers. The target population and units are fixed before treatment. With interference, use the complete assignment vector or a justified exposure mapping; individual no-interference notation alone cannot identify an economy-wide intervention. A difference of mean potential outcomes does not require the cross-world joint distribution. Conditioning on post-treatment migration, survival, enrollment or employment changes the estimand and needs separate justification.

\textbf{Statistical statements.} An estimator, confidence set, forecast or test specifies a sampling law, model class and loss/coverage criterion. Uniform \((1-\alpha)\)-coverage means \(\inf_{m\in\mathcal M}\Pr_m\{\tau(m)\in C_n\}\geq1-\alpha\), or an explicitly asymptotic version. The whole parameter space trivially has coverage, so informativeness requires a stated width, risk or power objective. A forecasting improvement is relative to a named benchmark, information set and evaluation population.

\textbf{Equilibrium and optimization.} An equilibrium comprises feasible plans, optimality, beliefs where needed, budget constraints and all market/resource clearing. The primitives or a stated benchmark define each correspondence \(\mathcal E(p;m)\). Nonemptiness is assumed or proved, never inferred just from compact policy sets. A selected-equilibrium target uses a declared selection rule; a robust target quantifies over all permitted equilibria. Use suprema or infima unless attainment is established. An \(\varepsilon\)-optimal policy is feasible and within \(\varepsilon\) of the finite optimum value. Report an empty feasible set separately.

\textbf{Welfare and accounting.} Normative utility, interpersonal normalization, social weights, numeraire, discounting and the use of public receipts are primitives. Behavioral utility need not coincide with the welfare criterion. Household welfare includes ownership income and rents once; adding profits again to compensating variation generally double-counts them. A monetary equivalent is defined by an explicit transfer and price convention, with monotonicity and range conditions. Average effects, mechanism contributions, transfers and welfare losses are different targets. Infinite sums require convergence and dynamic feasible plans require terminal or no-Ponzi conditions.

\textbf{Source versions.} A working-paper date, revision date and journal publication year are distinguished. A DOI for a journal version is not automatically the DOI of an earlier manuscript. Locators refer to the stated version and printed pages where specified. A metadata-only check does not verify a claim about a full-text conclusion. Every entry's source subsection narrows the provenance claim accordingly.
% END ECONOMICS STATEMENT
\end{document}
